% technique: doi-bien-cau-song
\textbf{Đổi biến cầu $v=r\,u$ cho sóng đối xứng cầu.}
\emph{$u=u(r,t)$ là nghiệm cần tìm} (biên độ tại khoảng cách $r$ tới tâm); \emph{$v=r\,u$ là hàm phụ} đặt ra để tính cho dễ.

\textbf{Vì sao đặt $v=ru$:} Laplacian của hàm chỉ phụ thuộc $r$ có số hạng phiền $\tfrac{2}{r}u_r$:
\[
  \Delta u = u_{rr}+\tfrac{2}{r}u_r .
\]
Nhân $u$ với $r$: $v_r=u+ru_r$, $v_{rr}=2u_r+ru_{rr}=r\big(u_{rr}+\tfrac{2}{r}u_r\big)=r\,\Delta u$. Nhân phương trình sóng với $r$:
\[
  r\,u_{tt}=c^2 r\,\Delta u+rf \;\Longrightarrow\; v_{tt}=c^2 v_{rr}+rf .
\]
Số hạng $\tfrac{2}{r}u_r$ \textbf{biến mất} $\Rightarrow$ $v$ thỏa \textbf{sóng 1D phẳng} trên $r>0$ (có nguồn $rf$ nếu có nguồn).

\textbf{Giải:} dùng d'Alembert cho $v$, rồi lấy lại $u=v/r$. Biên $v(0,t)=ru\big|_{r=0}=0$ (vì $u$ hữu hạn ở tâm) là \textbf{Dirichlet} $\Rightarrow$ thác triển dữ liệu/nguồn \textbf{lẻ} qua $r=0$ (xem ``Thác triển chẵn/lẻ'').
Dùng khi dữ liệu/nguồn đối xứng cầu (vd $\psi,f$ chỉ phụ thuộc $r$).
